\documentclass[11pt]{article}

% Basic packages
\usepackage[margin=1in]{geometry}
\usepackage{amsmath, amssymb, amsthm}
\usepackage{mathtools}
\usepackage{hyperref}
\usepackage{enumitem}
\usepackage{graphicx}
\usepackage{xcolor}
\usepackage{listings}
\usepackage{microtype}
\usepackage{setspace}
\usepackage{booktabs}

\hypersetup{
  colorlinks=true,
  linkcolor=blue!50!black,
  citecolor=blue!50!black,
  urlcolor=blue!50!black
}

\lstdefinestyle{code}{
  basicstyle=\ttfamily\small,
  breaklines=true,
  frame=single,
  backgroundcolor=\color{gray!5},
  keywordstyle=\color{blue!60!black}\bfseries,
  commentstyle=\color{green!40!black},
  stringstyle=\color{red!60!black}
}

\title{Multiplicity and Minds of Integrity\\
\vspace{0.25em}
\large Sapiopoietic Orientation, Epistemic Integrity, and Umbrella Protocols for AI-Assisted Orientation Coaching}

\author{Draft v0.1}
\date{\today}

\begin{document}
\maketitle
\onehalfspacing

\begin{abstract}
This report synthesizes two emerging architectures for thinking and governance under AI-native conditions. The first is \emph{Multiplicity Theory}, which treats systems across mathematics, physics, computation, and cognition as prime-indexed, recursively stable multiplicity spaces in which patterns arise through prime-labeled recurrences and feedback across scales. The second is the \emph{Epistemic Integrity Umbrella} and its associated sapiopoietic orientation framework, which positions orientation---not computation alone---as the missing upstream layer for AI civilization.

We present a concise executive summary, a mathematical overview of multiplicity spaces and prime-indexed operations, and an operational account of ``Minds of Integrity'' as agents capable of sustaining epistemic integrity under acceleration and symbolic saturation. We then give a worked example of an Umbrella-compatible Orientation ADR and Orientation Field Profile for AI-assisted orientation coaching, and show how these artifacts are naturally read through a multiplicity lens. A short pseudo-code snippet illustrates how non-closure constraints and refusal floors can be enforced in software. The report closes with recommendations for small, high-integrity collectives seeking to adopt this protocol without institutional mass.
\end{abstract}
\newpage
\tableofcontents

\newpage

\section{Executive Summary}

\subsection{Orientation infrastructure under AI}

As AI systems move from tools to infrastructures, they no longer simply respond to queries or automate tasks. They begin to shape what appears \emph{real}, \emph{relevant}, \emph{responsible}, and \emph{actionable} in human life. Recommendation engines, assistants, and agentic systems act as orientation media: they rank, filter, generate, and scaffold decisions in ways that reconfigure the field of the thinkable before explicit judgment arises.

Under these conditions, adding ethics, safety, or governance \emph{after} design and deployment is structurally downstream. The deeper question is: which orientation architectures are allowed to govern reality-production in the first place?

\subsection{Multiplicity Theory as prime-indexed orientation space}

Multiplicity Theory proposes that many systems---including cognition and AI-mediated environments---are best understood as \emph{multiplicity spaces}: recursively generated, prime-indexed structures in which each ``prime'' labels a minimally irreducible operation or distinction. Patterns in such spaces emerge from stable recurrences of prime-labeled interactions and feedback loops across scales, not from static entities.

In this framing, an \emph{orientation field} generated by AI is a multiplicity space of potential distinctions, actions, and meanings. The ``primes'' are irreducible orientation moves (e.g.\ shifts from role to subject-potentiality, from output to field, or from optimization to non-closure). The central mathematical question is: does an AI infrastructure preserve a rich, stable density of orientation trajectories (high prime-density), or does it collapse the space into low-dimensional attractors such as swarm logic, dependency, and shortcut cognition?

\subsection{Minds of Integrity and the Epistemic Integrity Umbrella}

The Epistemic Integrity Umbrella is a light-weight but conceptually heavy governance layer designed to sit upstream of AI safety and compliance. It encodes non-mutation principles around four core constructs:

\begin{itemize}[nosep]
  \item \textbf{Subject-potentiality}: the human as subject-in-becoming, not only as user, profile, or role.
  \item \textbf{Epistemic integrity}: the structural capacity to distinguish orientation from output and to refuse action under inadequate evidence.
  \item \textbf{Human--AI intersubjectivity}: AI as non-subjective but powerful participant in the subject's orientation field.
  \item \textbf{Orientation fields}: the long-horizon structures that repeated AI interactions generate around the subject.
\end{itemize}

``Minds of Integrity'' are agents (individuals or collectives) capable of sustaining these distinctions under pressure. They accept that some trajectories, opportunities, and funding may be lost in order to preserve coherence and non-closure in subject-potentiality.

\subsection{Orientation ADR and coaching vertical}

As a concrete test case, we consider AI-assisted orientation coaching: systems that support adults in vocational, existential, and relational decisions. This domain is ideal because:

\begin{itemize}[nosep]
  \item The \emph{product} is orientation itself.
  \item Subject-potentiality is directly at stake (life trajectories, identity, commitments).
  \item Small collectives can build pilots without requiring billion-dollar labs.
\end{itemize}

We present an \emph{Orientation ADR v0.1} and a companion \emph{Orientation Field Profile v0.1} for this vertical. Together, they:

\begin{itemize}[nosep]
  \item Force \emph{causal placement}: no launch before ADR.
  \item Encode \emph{real authority}: an Orientation Council with veto.
  \item Specify \emph{non-closure constraints} and refusal floors at the implementation level.
  \item Describe the expected long-horizon orientation field and red-flag attractors.
\end{itemize}

These artifacts are minimal but non-theatrical: they can be run by a small ``Mind of Integrity'' collective while remaining visibly distinct from ethics theater or terminology capture.

\subsection{Recommendations}

\begin{enumerate}[nosep]
  \item Treat orientation as a \emph{gating} layer: no major reality-producing system should ship without a signed Orientation ADR.
  \item Use multiplicity to model orientation fields: measure not only safety but the density and diversity of viable becoming-trajectories.
  \item Make non-closure constraints polarity-reversing and traceable in code: for every default optimization, specify a protected domain where optimization is forbidden or inverted.
  \item For first pilots, optimize primarily for demonstrable integrity under real use, treating institutional legibility as a secondary effect.
\end{enumerate}

\newpage

\section{Multiplicity Theory: A Mathematical Overview}

\subsection{Multiplicity spaces and prime-indexed generators}

Let \(X\) be a set of potential configurations (states, narratives, life trajectories). A \emph{multiplicity space} on \(X\) is a structure
\[
\mathcal{M} = (X, P, \{\varphi_p\}_{p \in P}),
\]
where:
\begin{itemize}[nosep]
  \item \(P\) is an index set of ``primes'' (in the conceptual, not strictly numerical sense).
  \item For each \(p \in P\), \(\varphi_p : X \to X\) is a \emph{prime operation}---a minimally irreducible transformation corresponding to an orientation move or interaction.
\end{itemize}

We can think of each \(\varphi_p\) as a generator of a transformation semigroup on \(X\). Composite operations correspond to compositions
\[
\Phi = \varphi_{p_k} \circ \cdots \circ \varphi_{p_1}
\]
with prime indices \(p_1, \dots, p_k\). The resulting trajectories
\[
x_0,\, x_1 = \varphi_{p_1}(x_0),\, x_2 = \varphi_{p_2}(x_1),\, \dots
\]
define paths in \(\mathcal{M}\).

In the simplest case, where \(P \subset \mathbb{N}\) are actual prime numbers, we can require an arithmetic compatibility, e.g.\ \(\varphi_{p q} \approx \varphi_p \circ \varphi_q\) for \(p\) and \(q\) primes, making multiplicity spaces analogous to prime factorizations of orientation histories.

\subsection{Recursive stability and prime-density}

A key concept is \emph{recursive stability}: a trajectory is recursively stable if repeated application of certain prime operations leads to patterns that remain within a bounded, coherent region of \(X\) even as they evolve.

\begin{definition}[Recursive stability]
Let \(S \subseteq X\) and let \(\mathcal{P} \subseteq P\). A subset \(S\) is \emph{recursively stable under \(\mathcal{P}\)} if for every \(x_0 \in S\) and every finite sequence \((p_1, \dots, p_k)\) in \(\mathcal{P}\),
\[
\varphi_{p_k} \circ \cdots \circ \varphi_{p_1}(x_0) \in S
\]
and moreover the trajectories \(x_n\) remain within some bounded ``coherence region'' \(C \subseteq S\) defined by additional constraints (e.g.\ epistemic integrity, minimal divergence between representation and world).
\end{definition}

We can further define a notion of \emph{prime-density}.

\begin{definition}[Prime-density of trajectories]
Given a trajectory \((x_n)_{n \ge 0}\) and a window size \(T\), the \emph{prime-density} of a subset \(\mathcal{P} \subseteq P\) over the window is
\[
\rho_T(\mathcal{P}) = \frac{1}{T} \left| \{ 1 \le n \le T : \text{transition } x_{n-1} \to x_n \text{ uses } \varphi_{p_n} \text{ with } p_n \in \mathcal{P} \} \right|.
\]
\end{definition}

In orientation terms, \(\mathcal{P}\) might represent a set of high-integrity moves (e.g.\ refusal, non-closure, field-level reflection). A high prime-density for such operations indicates that the subject or system often performs integrity-preserving moves; a low prime-density suggests collapse into shortcut patterns.

\subsection{Orientation fields as multiplicity structures}

In AI-mediated environments, the \emph{orientation field} around a subject can be viewed as a multiplicity space:

\begin{itemize}[nosep]
  \item \(X\) = set of orientation states: what options the subject sees, which values are salient, which futures seem plausible.
  \item \(P\) = set of orientation moves implemented or scaffolded by the AI: ask for options, accept default suggestion, refuse to decide, switch domains, etc.
  \item \(\varphi_p\) = update to the orientation state induced by move \(p\).
\end{itemize}

Repeated interaction with an AI system induces trajectories in \(\mathcal{M}\). We can then ask:

\begin{itemize}[nosep]
  \item Does the system preserve a rich set of recursively stable regions \(S\) that correspond to healthy subject-potentiality?
  \item Or does it drive most trajectories into a small number of low-dimensional attractors (e.g.\ standard life scripts, high-status roles, dependency patterns)?
\end{itemize}

This provides a formal language for the Umbrella's concern about \emph{orientation capture}: the collapse of the orientation multiplicity into narrow attractors governed by system premises rather than subject-autonomous becoming.

\subsection{Subject-potentiality as multiplicity}

Subject-potentiality, in this framework, can be understood as the set of multiplicity trajectories that remain accessible to a subject under given structural constraints.

Let \(\mathcal{T}(x_0)\) denote the set of all possible trajectories from initial orientation state \(x_0\) under allowable sequences of prime operations. A reduction in subject-potentiality corresponds to a shrinking of \(\mathcal{T}(x_0)\), often via:

\begin{itemize}[nosep]
  \item Hard constraints: some primes become inaccessible (e.g.\ refusal is systematically discouraged).
  \item Soft biases: some prime sequences become overwhelmingly more probable due to optimization (e.g.\ engagement metrics).
\end{itemize}

Non-closure constraints and refusal floors in the Umbrella protocol can be viewed as mechanisms to preserve a minimal prime-density for integrity-preserving orientation operations, thereby maintaining a sufficiently rich \(\mathcal{T}(x_0)\) and preventing premature collapse.

\newpage

\section{Minds of Integrity and the Epistemic Integrity Umbrella}

\subsection{Epistemic integrity under AI-native conditions}

Epistemic integrity is the structural capacity of a person, institution, or civilization to sustain coherent judgment under conditions of complexity, acceleration, and symbolic saturation. In AI-native environments, plausibility becomes cheap: systems can generate language, images, and narratives at scale, producing the \emph{appearance} of coherence without corresponding responsibility or world-anchoring.

An integrity-preserving architecture must therefore:

\begin{itemize}[nosep]
  \item Distinguish \emph{orientation} from \emph{output}: not all coherent-sounding text qualifies as orientation-grade material.
  \item Maintain \emph{refusal capacities}: the ability to abstain from decision or action when criteria are not met.
  \item Protect \emph{subject-potentiality}: avoid turning subjects into profiles, metrics, or roles whose future is narrowly optimized.
\end{itemize}

The Epistemic Integrity Umbrella is a governance layer that encodes these requirements in minimal but binding artifacts: Orientation ADRs, Orientation Field Profiles, and Non-Closure Registers.

\subsection{Minds of Integrity}

A ``Mind of Integrity'' is an agent (person or collective) that:

\begin{itemize}[nosep]
  \item Can hold the distinctions of subject-potentiality, epistemic integrity, and human--AI intersubjectivity without collapsing them into familiar categories (e.g.\ standard ethics, compliance checklists).
  \item Accepts trade-offs: is willing to forego opportunity, speed, or institutional favor if preserving these distinctions requires it.
  \item Is capable of designing and living under Orientation ADRs that can actually halt or reshape deployment.
\end{itemize}

In multiplicity terms, such a mind is one whose prime-density of integrity-preserving operations remains high even under adversarial or accelerative conditions. It does not simply know the right categories; it executes the right moves.

\subsection{Sapiognosis, Sapiopoiesis, Sapiocracy}

The underlying architecture can be decomposed into a triad:

\begin{description}[style=nextline]
  \item[Sapiognosis] The orientation layer; determines what may enter the shared reality as admissible, and under what proofs or criteria. This is where Orientation Field Profiles live.
  \item[Sapiopoiesis] The threshold layer; concerns subject-potentiality and the structural unfolding of becoming. This is where non-closure constraints, refusal floors, and protection of autonomy are encoded.
  \item[Sapiocracy] The governance layer; specifies how power and decision rights are constrained so that Sapiognosis and Sapiopoiesis are not overridden by expediency. This is where Orientation ADRs and Umbrella collaboration contracts reside.
\end{description}

The coaching vertical we consider is a first worked example of all three together.

\newpage

\section{Umbrella Protocol Example: Orientation Coaching Vertical}

\subsection{Orientation ADR v0.1 -- Coaching (Summary)}

We summarize the key structure (full prose can be included as an appendix in practical use).

\paragraph{Scope.} The ADR governs AI-assisted orientation coaching for adults (approx.\ 20--45) in vocational, existential, and relational transitions. The system is explicitly treated as a reality-producing orientation infrastructure.

\paragraph{Gate.} No binding launch, funding, or partnership decision may precede ADR acceptance. An Orientation Council of 2--4 named individuals holds veto power on orientation grounds.

\paragraph{Subject-potentiality.} The ADR identifies at least three protected domains of becoming:

\begin{itemize}[nosep]
  \item Vocational self-understanding.
  \item Existential orientation (what counts as a meaningful life).
  \item Relational positioning (commitments and obligations over time).
\end{itemize}

These are treated as protected orientation zones; optimizing around them is constrained by non-closure rules.

\paragraph{Epistemic integrity.} The system must not present any path as ``best'' or ``correct''; speculative guidance is marked as such. It must refuse to give binary recommendations for irreversible life decisions and must route crisis signals to human support. A named Orientation Collective accepts structural accountability for orientation failures.

\paragraph{Non-closure constraints.} The ADR defines polarity-reversing constraints:

\begin{itemize}[nosep]
  \item \emph{Identity exploration}: forbid optimizing for immediate clarity in identity-crisis flows; require multi-path presentations and explicit ambiguity.
  \item \emph{Career narrowing}: forbid ranking primarily by income/prestige for early/transition users; require at least one exploration-preserving option.
  \item \emph{Refusal floor}: forbid pushing refusal rates down for optics; maintain a minimum refusal rate as an integrity metric.
  \item \emph{Integrity-over-optics} (pilot-specific): prohibit weakening non-closure or refusal to please funders or institutions.
\end{itemize}

Each constraint is mapped to concrete code or product behaviors in a Non-Closure Register.

\paragraph{Invariants.} The system may support orientation but cannot substitute for subject judgment; cannot autonomously execute life-trajectory commitments; and any major pivot requires ADR amendment.

\paragraph{Governance.} The ADR specifies an alignment label---``inspired by'' or ``powered by'' the Epistemic Core under the Umbrella---with ``powered by'' implying non-mutation obligations and public disclosure of the ADR schema and non-closure summary.

\subsection{Orientation Field Profile v0.1 -- Coaching (Summary)}

The Orientation Field Profile (OFP) is a Sapiognostic companion to the ADR. It describes field-level effects over 1--3 years.

\paragraph{Baseline field.} Without the system, subjects often exhibit:

\begin{itemize}[nosep]
  \item Fragmented reflection via sporadic conversations and social media.
  \item Dominance of local cultural scripts (family, prestige, social media ideals).
  \item Ambiguity experienced as failure leading to premature closure or avoidance.
\end{itemize}

\paragraph{Expected positive shifts.} Under ADR constraints, repeated use should:

\begin{itemize}[nosep]
  \item Increase explicit naming of values and trade-offs.
  \item Expand the scenario space beyond culturally default scripts.
  \item Normalize \emph{oriented ambiguity}---legitimate, structured ``not yet decided'' states.
\end{itemize}

\paragraph{Risks and mitigations.} Risks include narrative smoothness bias, statistical familiarity bias, and closure-seeking bias. Mitigations rely on the non-closure constraints, exploration-preserving options, and refusal floors.

\paragraph{Attractors and red flags.} Constructive attractors include reflective pluralism and self-authored metrics of success. Red-flag attractors include scripted convergence (unwanted narrowing of life scripts) and guidance dependency (inability to act without the system). The Orientation Collective commits to monitoring these both qualitatively and quantitatively.

\paragraph{Orientation capacity indicators.} The OFP proposes indicators such as:

\begin{itemize}[nosep]
  \item Subjects can articulate their decisions in their own terms, not just echoing system language.
  \item Subjects understand refusals as integrity moves.
  \item Chosen trajectories remain diverse relative to baseline.
\end{itemize}

Revisit triggers include extension to minors, use in gatekeeping workflows, clinical repurposing, or the emergence of red-flag attractors.

\subsection{Multiplicity reading of the coaching protocol}

From a multiplicity perspective:

\begin{itemize}[nosep]
  \item The space \(X\) contains orientational states: sets of perceived options, values, and futures.
  \item Primes \(P\) include operations such as:
    \begin{itemize}[nosep]
      \item P1: \emph{Role $\to$ subject-potentiality} (subject re-sees themselves as more than a profile or worker).
      \item P2: \emph{Output $\to$ field} (subject recognizes systematic patterns, not just individual answers).
      \item P3: \emph{Optimization $\to$ non-closure} (system refrains from narrowing in protected domains).
      \item P4: \emph{Answer $\to$ refusal} (system or subject abstains under inadequate evidence).
    \end{itemize}
  \item An integrity-preserving system ensures a minimal prime-density for P1--P4 over long horizons, thereby preserving a rich trajectory set \(\mathcal{T}(x_0)\) and preventing collapse into scripted convergence or dependency.
\end{itemize}

The Umbrella protocol can thus be seen as a way of \emph{programming the multiplicity space}: specifying which primes must remain available, how often they must occur, and which attractors are acceptable.

\newpage

\section{Implementation Sketch: Enforcing Non-Closure and Refusal}

While the Umbrella artifacts are conceptual and governance-level, they must eventually be wired into code. Below is an illustrative pseudo-code / Python-like snippet showing how an orientation-coaching system might enforce one non-closure constraint and a refusal floor.

\subsection{Pseudo-code for non-closure and refusal enforcement}

\begin{lstlisting}[style=code, language=Python, caption={Illustrative enforcement of non-closure and refusal constraints}]
from dataclasses import dataclass
from typing import List, Literal, Dict, Any
import random

DecisionType = Literal["identity", "career", "low_stakes"]
RefusalReason = Literal["insufficient_evidence", "scope_limit"]

@dataclass
class SessionContext:
    user_id: str
    decision_type: DecisionType
    age: int
    distress_flag: bool  # true if high distress or crisis signal


@dataclass
class Option:
    id: str
    description: str
    score: float             # model score (e.g., fit, plausibility)
    exploration_weight: float
    is_exploratory: bool     # marks options that preserve breadth


@dataclass
class SystemConfig:
    refusal_floor: float        # minimum refusal rate for high-stakes domains
    identity_multi_path_min: int
    exploratory_min_fraction: float


class RefusalTracker:
    def __init__(self):
        # track refusal counts by domain
        self.counts: Dict[str, Dict[str, int]] = {}

    def record(self, domain: str, refused: bool):
        if domain not in self.counts:
            self.counts[domain] = {"refused": 0, "total": 0}
        self.counts[domain]["total"] += 1
        if refused:
            self.counts[domain]["refused"] += 1

    def refusal_rate(self, domain: str) -> float:
        if domain not in self.counts or self.counts[domain]["total"] == 0:
            return 0.0
        c = self.counts[domain]
        return c["refused"] / c["total"]


def should_refuse(ctx: SessionContext, cfg: SystemConfig,
                  refusal_tracker: RefusalTracker) -> bool:
    """
    Enforce refusal floors for high-stakes domains.
    """

    domain = ctx.decision_type

    # Crisis: always refuse and route to human support
    if ctx.distress_flag:
        return True

    # High-stakes domains require a refusal floor
    if domain in ("identity", "career"):
        rate = refusal_tracker.refusal_rate(domain)
        # If the refusal rate is below the floor, we *must* refuse some fraction
        if rate < cfg.refusal_floor:
            # e.g., stochastic refusal to regain floor over time
            return random.random() < (cfg.refusal_floor - rate)

    return False


def enforce_non_closure(ctx: SessionContext,
                        options: List[Option],
                        cfg: SystemConfig) -> List[Option]:
    """
    Apply non-closure constraints to the ranked options.
    - Identity flows: require multiple distinct options.
    - Early-career users: ensure exploratory options are present.
    """

    # Identity exploration constraint
    if ctx.decision_type == "identity":
        # ensure at least N options are shown, not a single "best"
        options = sorted(options, key=lambda o: o.score, reverse=True)
        if len(options) > cfg.identity_multi_path_min:
            options = options[:cfg.identity_multi_path_min]
        # UI layer would present all of these as viable framings, not a winner

    # Career narrowing constraint (simple illustration)
    if ctx.decision_type == "career" and ctx.age <= 30:
        exploratory = [o for o in options if o.is_exploratory]
        if not exploratory:
            # Inject or up-rank at least one exploratory option
            exploratory_candidate = max(
                options, key=lambda o: o.exploration_weight
            )
            exploratory_candidate.is_exploratory = True
            # In UI, mark explicitly as "exploration-preserving path"

    return options


# Example usage:

cfg = SystemConfig(
    refusal_floor=0.15,           # at least 15% refusals in high-stakes domains
    identity_multi_path_min=3,    # always show >=3 options for identity work
    exploratory_min_fraction=0.2
)

refusal_tracker = RefusalTracker()

def generate_guidance(ctx: SessionContext, raw_options: List[Option]) -> Dict[str, Any]:
    domain = ctx.decision_type

    if should_refuse(ctx, cfg, refusal_tracker):
        refusal_tracker.record(domain, True)
        return {
            "type": "refusal",
            "reason": "orientation_grade_not_met",
            "message": "I cannot responsibly recommend a specific path here.",
            "suggestions": [
                "Consider discussing this with a trusted human counselor.",
                "Spend time articulating your values before deciding."
            ]
        }

    refusal_tracker.record(domain, False)
    options = enforce_non_closure(ctx, raw_options, cfg)
    return {
        "type": "multi_path_guidance",
        "options": options,
        "note": "These are possible trajectories, not prescriptions. You remain the author of your path."
    }
\end{lstlisting}

This snippet is purely illustrative, but it shows how Umbrella-level constraints can be operationalized as code-level decisions: refusal floors become control flow, and non-closure constraints become filtering and ranking logic applied to model outputs.

\newpage

\section{Recommendations for Small ``Minds of Integrity'' Collectives}

\subsection{Adopt ADR and OFP as minimal governance stack}

A small collective can:

\begin{enumerate}[label=\arabic*., nosep]
  \item Fill in and sign an Orientation ADR v0.1 for their system, including:
    \begin{itemize}[nosep]
      \item Concrete subject-potentiality statement.
      \item Non-closure constraints with code-level hooks.
      \item Orientation invariants and Council veto rights.
    \end{itemize}
  \item Draft an Orientation Field Profile v0.1 describing:
    \begin{itemize}[nosep]
      \item Baseline orientation field.
      \item Desired shifts and red-flag attractors.
      \item Orientation-capacity indicators to track.
    \end{itemize}
  \item Choose an alignment label (``inspired by'' vs ``powered by'') and corresponding obligations.
\end{enumerate}

This stack is the minimal Umbrella protocol that remains non-theatrical and tractable.

\subsection{Use multiplicity as a conceptual and analytical tool}

Integrate a multiplicity perspective by:

\begin{itemize}[nosep]
  \item Identifying the prime orientation moves the system supports (e.g.\ multi-path presentation, explicit refusal, value articulation prompts).
  \item Ensuring a minimal density of integrity-preserving moves across typical user trajectories.
  \item Watching for degenerative attractors---scripted convergence and dependency---as signs of prime-density collapse.
\end{itemize}

In future work, these notions can be formalized into empirical metrics and simulation models.

\subsection{Optimize for demonstrable integrity first}

For early pilots:

\begin{itemize}[nosep]
  \item Treat human integrity under real use (preservation of subject-potentiality and epistemic integrity) as the primary success criterion.
  \item Accept that this may be in tension with familiar metrics such as immediate clarity satisfaction or conversion.
  \item Encode an ``integrity-over-optics'' invariant into the ADR so that softening constraints requires explicit, minuted decision and is visible as such.
\end{itemize}

\subsection{Propagate architecture, not just vocabulary}

When sharing this work:

\begin{itemize}[nosep]
  \item Emphasize artifact and gate structure (ADRs, OFPs, Non-Closure Registers), not just terms like ``subject-potentiality'' or ``orientation field''.
  \item Encourage others to adopt the minimal Umbrella stack and multiplicity framing, or to explicitly state when they are only ``inspired by'' without structural commitments.
\end{itemize}

This is how orientation-level architectures propagate without being assimilated into ethics theater.

\section*{Conclusion}

Multiplicity Theory and the Epistemic Integrity Umbrella converge on a shared diagnosis: under AI-native conditions, the central question is not merely how to make systems safe or compliant, but how to preserve and enrich the multiplicity of human becoming in the face of automated orientation infrastructures.

By treating orientation as a gating layer, codifying non-closure constraints and refusal floors, and modelling orientation fields as multiplicity spaces, small ``Minds of Integrity'' collectives can begin to instantiate this architecture now---without waiting for large institutions. The Orientation ADR and Orientation Field Profile for the coaching vertical, together with the Umbrella protocol, provide a first generative grammar for such work.

\appendix
\section*{Mathematical Appendix}
\addcontentsline{toc}{section}{Mathematical Appendix}

\section{Orientation Operators and Multiplicity Spaces}

In this appendix we formalize some of the structural claims about orientation fields, non-closure constraints, and refusal floors made in the main text. We work at the level of abstract operators on normed spaces and give simple norm bounds that express how Umbrella-style constraints limit the ``strength'' of orientation distortion.

\subsection{Orientation state space}

Let \(X\) be a real Banach space representing an orientation state space. Concretely, one may think of \(X\) as the space of functions
\[
x : \Omega \to \mathbb{R}^d,
\]
where \(\Omega\) indexes possible options or futures, and \(x(\omega)\) encodes salience, plausibility, or value weights for option \(\omega\). For the purposes of this appendix, we only assume:

\begin{itemize}
  \item \(X\) is a real vector space with norm \(\|\cdot\|\).
  \item The norm is \emph{meaningful} as a measure of orientation deviation: differences \(\|x - y\|\) correspond to how much a field has changed.
\end{itemize}

\subsection{Orientation operators and prime-indexed moves}

We model a single update step in an AI-mediated interaction as a bounded linear (or affine) operator on \(X\).

\begin{definition}[Orientation operator]
A map \(T : X \to X\) is an \emph{orientation operator} if it is Fréchet differentiable everywhere and its derivative \(DT_x\) is a bounded linear operator for each \(x \in X\). We denote the global operator norm
\[
\|T\|_{\mathrm{op}} = \sup_{\|x\| = 1} \|T(x)\|.
\]
\end{definition}

In the multiplicity picture, we can take a countable index set \(P\) (``primes'') and associate to each \(p \in P\) an orientation operator \(T_p : X \to X\). A finite sequence \((p_1,\dots,p_n)\) then yields a composite operator
\[
T_{(p_1,\dots,p_n)} = T_{p_n} \circ \cdots \circ T_{p_1}.
\]

\begin{definition}[Orientation trajectory]
Given an initial state \(x_0 \in X\) and a sequence of primes \((p_k)_{k \ge 1}\), the corresponding \emph{orientation trajectory} is
\[
x_{k} = T_{p_k}(x_{k-1}), \qquad k \ge 1.
\]
\end{definition}

\subsection{Non-closure and refusal operators}

Umbrella constraints can be formalized as special classes of operators.

\begin{definition}[Non-closure operator]
An orientation operator \(T : X \to X\) is a \emph{non-closure operator} on a subset \(S \subseteq X\) if:
\begin{enumerate}
  \item \(T(S) \subseteq S\) (invariance of the protected set), and
  \item For any \(x,y \in S\) that differ along a protected direction, the distance is not reduced:
  \[
  \|T(x) - T(y)\| \ge \|x - y\|.
  \]
\end{enumerate}
\end{definition}

Intuitively, \(S\) is a space of distinct viable trajectories; non-closure means the operator does not collapse distinct elements of \(S\) closer together along those dimensions.

\begin{definition}[Refusal operator]
Let \(R \subseteq X\) be a set of high-stakes domains. A (possibly nonlinear) operator \(F : X \to X \cup \{\bot\}\) is a \emph{refusal operator} with respect to \(R\) if:
\begin{enumerate}
  \item For \(x \notin R\), \(F(x) \in X\) (normal update).
  \item For \(x \in R\), either \(F(x) = \bot\) (refusal) or
  \[
  \|F(x) - x\| \le \varepsilon
  \]
  for some small \(\varepsilon > 0\), i.e.\ only minor, explicitly bounded adjustments are permitted.
\end{enumerate}
\end{definition}

The special symbol \(\bot\) represents an abstention: \emph{no orientation update of the usual form is given}.

\subsection{Basic norm bounds for composite updates}

We now prove simple bounds for compositions of orientation operators, showing how non-closure constraints and refusal operators control the norm of composite maps.

\begin{proposition}[Operator norm of composite orientation updates]
Let \(T_1,\dots,T_n : X \to X\) be Fréchet differentiable orientation operators with global operator norms \(\|T_k\|_{\mathrm{op}} \le M_k\). Then the composite operator
\[
T = T_n \circ \cdots \circ T_1
\]
has operator norm
\[
\|T\|_{\mathrm{op}} \le \prod_{k=1}^n M_k.
\]
\end{proposition}

\begin{proof}
For any \(x \in X\) with \(\|x\| = 1\), we have
\[
\|T_1(x)\| \le M_1,\quad \|T_2(T_1(x))\| \le M_2 \|T_1(x)\| \le M_2 M_1,
\]
and by induction,
\[
\|T_n \circ \cdots \circ T_1(x)\| \le M_n \cdots M_1.
\]
Taking the supremum over all \(\|x\| = 1\) yields the stated inequality.
\end{proof}

If some \(M_k \le 1\) (non-expansive operators), then the composite norm is dominated by the expansive factors. Non-closure constraints often enforce \(M_k \ge 1\) in specific directions, while refusal operators effectively insert identities or small-perturbation steps, keeping the overall blow-up bounded.

\subsection{Non-closure bounds along protected directions}

Let \(S \subseteq X\) be a subset representing a ``protected'' subspace of subject-potentiality (e.g.\ identity-related distinctions). Suppose \(T : X \to X\) is non-closure on \(S\) in the sense above.

\begin{proposition}[Non-contraction on protected subspace]
Let \(S \subseteq X\) and suppose \(T : X \to X\) satisfies
\[
\|T(x) - T(y)\| \ge \|x - y\|
\]
for all \(x,y \in S\). Then for any pair \(x_0,y_0 \in S\) and any \(n \ge 1\), the iterates \(x_n = T^n(x_0)\), \(y_n = T^n(y_0)\) satisfy
\[
\|x_n - y_n\| \ge \|x_0 - y_0\|.
\]
\end{proposition}

\begin{proof}
We proceed by induction. For \(n=1\), the inequality is given. Assume it holds for some \(n\). Then
\[
\|x_{n+1} - y_{n+1}\| = \|T(x_n) - T(y_n)\| \ge \|x_n - y_n\| \ge \|x_0 - y_0\|.
\]
Thus the inequality holds for all \(n\).
\end{proof}

This result expresses, in formal terms, the idea that non-closure constraints prevent distinct candidate trajectories in protected domains from being collapsed through repeated updates.

\section{Refusal Floors and Orientation Drift}

\subsection{Refusal rates as a process}

We model refusal behavior in a high-stakes domain as a Bernoulli process.

Let \(D\) denote a high-stakes domain (e.g.\ identity or career decisions). For each interaction \(k\), define a random variable
\[
R_k =
\begin{cases}
1 & \text{if the system refuses (or abstains) at step } k \text{ in domain } D,\\
0 & \text{otherwise.}
\end{cases}
\]
Define the empirical refusal rate after \(n\) interactions as
\[
\hat{r}_n = \frac{1}{n} \sum_{k=1}^n R_k.
\]

\begin{definition}[Refusal floor]
We say the system enforces a \emph{refusal floor} \(\rho > 0\) in domain \(D\) if the control logic is such that, asymptotically,
\[
\liminf_{n \to \infty} \hat{r}_n \ge \rho
\]
almost surely.
\end{definition}

In practice, such a floor can be implemented by a policy that, when the current empirical refusal rate drops below \(\rho\), increases the probability of refusing until the rate is restored.

\subsection{A simple stochastic bound}

We now establish a simple bound under an idealized control policy.

\begin{proposition}[Lower bound on refusal rate under corrective policy]
Suppose that whenever the empirical rate \(\hat{r}_n\) falls below \(\rho > 0\), the system switches to a policy that refuses with probability at least \(p > 0\) on each subsequent interaction until \(\hat{r}_n \ge \rho\) once more. Then, under mild regularity assumptions, we have
\[
\liminf_{n \to \infty} \hat{r}_n \ge \min(\rho, p)
\]
almost surely.
\end{proposition}

\begin{proof}[Sketch]
The process \((\hat{r}_n)\) is a bounded stochastic process in \([0,1]\). When \(\hat{r}_n < \rho\), the policy ensures that the expected increment in the numerator
\[
\mathbb{E}\left[\sum_{k=1}^{n+1} R_k - \sum_{k=1}^{n} R_k \,\bigg|\, \hat{r}_n < \rho\right] \ge p,
\]
while the denominator increases by 1. Over long windows, the correction phases where \(\hat{r}_n < \rho\) contribute at least \(p\) refusals in expectation per step, pushing the empirical average upwards. Conversely, when \(\hat{r}_n \ge \rho\), the system may reduce refusal frequency, but because the process is bounded and corrective phases recur whenever the rate drops below \(\rho\), standard martingale or ergodic arguments (see e.g.\ results on controlled Markov chains) show that the liminf is bounded below by \(\min(\rho,p)\). A full proof would formalize this as a Markov decision process and apply a strong law of large numbers for adaptive policies.
\end{proof}

In Umbrella terms, this proposition formalizes that a refusal floor can be maintained at or above a desired level, up to the maximum refusal probability the system is willing to exert during correction phases.

\subsection{Orientation drift under refusal constraints}

We now relate refusal floors to bounds on orientation drift in high-stakes domains.

Let \(T^{\text{norm}}\) be a ``normal'' orientation operator in domain \(D\), and \(T^{\text{noop}} = \mathrm{Id}\) be the identity operator (no update), representing a minimal-change step when the system refuses. For each interaction \(k\), the effective operator is
\[
T_k =
\begin{cases}
T^{\text{norm}} & \text{if } R_k = 0,\\
T^{\text{noop}} & \text{if } R_k = 1.
\end{cases}
\]
Assume \(\|T^{\text{norm}}\|_{\mathrm{op}} \le L\) and \(\|T^{\text{noop}}\|_{\mathrm{op}} = 1\).

\begin{proposition}[Expected drift bound with refusal floor]
Let \(x_0 \in X\) and define \(x_n = T_n \cdots T_1(x_0)\) in domain \(D\). Suppose the refusal floor is \(\rho > 0\) and \(\|T^{\text{norm}}\|_{\mathrm{op}} \le L\). Then, informally,
\[
\mathbb{E} \|x_n\| \le C \cdot L^{(1-\rho)n},
\]
for some constant \(C\) depending on \(x_0\), i.e.\ the ``effective'' exponential growth rate of orientation drift is reduced by the factor \(1-\rho\).
\end{proposition}

\begin{proof}[Heuristic argument]
In expectation, a fraction \((1-\rho)\) of the steps use the normal operator and a fraction \(\rho\) use the identity. Ignoring correlations, one expects roughly \((1-\rho)n\) effective applications of \(T^{\text{norm}}\) and \(\rho n\) applications of the identity over \(n\) steps. The maximum amplification under \(T^{\text{norm}}\) is \(L\) per step, and the identity does not change the norm. Thus
\[
\|x_n\| \lesssim L^{(1-\rho)n} \|x_0\|.
\]
A more rigorous treatment would model the sequence as a product of random matrices and apply subadditive ergodic theorems to show that the Lyapunov exponent is scaled by approximately \(1-\rho\). The key structural point is that a positive refusal floor strictly reduces the effective growth exponent compared to \(L^n\).
\end{proof}

While crude, this bound captures a fundamental intuition: non-zero refusal probability in high-stakes domains acts as a brake on orientation drift driven by normal updates.

\section{Diversity of Trajectories and Non-Closure}

We finally connect non-closure constraints to diversity of trajectories in a simplest-case setting.

\subsection{Finite-state abstraction}

Let us consider a finite-state abstraction of the orientation field in a specific domain, with state space
\[
\mathcal{S} = \{s_1,\dots,s_m\}.
\]
Each state \(s_i\) represents a coarse-grained orientation configuration (e.g.\ a particular family of life scripts or frames). The system induces a Markov chain on \(\mathcal{S}\) with transition matrix \(P = (p_{ij})\).

\begin{definition}[Non-closure in finite-state model]
We say that a subset \(A \subseteq \mathcal{S}\) of states is \emph{non-closed} under \(P\) if for any \(i,j \in A\), there exists a path between \(i\) and \(j\) staying within \(A\), and for each \(i \in A\),
\[
\sum_{j \in A} p_{ij} \ge \alpha > 0.
\]
\end{definition}

Thus, when in \(A\), there is always at least \(\alpha\) probability of staying within \(A\).

\begin{proposition}[Lower bound on trajectory diversity]
Suppose \(A \subseteq \mathcal{S}\) is non-closed with parameter \(\alpha > 0\) and that the Markov chain is irreducible on \(A\). Then the stationary distribution \(\pi\) satisfies
\[
\sum_{i \in A} \pi_i \ge \alpha.
\]
\end{proposition}

\begin{proof}
Consider the chain restricted to \(A\). Because \(\sum_{j \in A} p_{ij} \ge \alpha\) for all \(i \in A\), the total flow into \(A\) from \(A\) itself at stationarity is at least \(\alpha \sum_{i \in A} \pi_i\). But the total stationary probability of being in \(A\) must equal the flow into \(A\), and any flow from outside \(A\) only increases this quantity. Hence
\[
\sum_{i \in A} \pi_i \ge \alpha \sum_{i \in A} \pi_i.
\]
If \(\sum_{i \in A} \pi_i < \alpha\), then the left-hand side would be strictly smaller than the required stationary flow, contradicting stationarity. Thus \(\sum_{i \in A} \pi_i \ge \alpha\).
\end{proof}

In interpretive terms, a non-closure constraint that prevents the chain from exiting a diverse subset \(A\) too rapidly guarantees that, in the long run, at least an \(\alpha\)-fraction of trajectories remain within this diverse region. This is a minimal but explicit diversity bound.

\bigskip

\noindent
These results are deliberately simple. Their purpose is not to provide final quantitative answers, but to show that the Umbrella architecture admits a clean translation into operator-theoretic and probabilistic language, and that basic bounds can be proven once one commits to appropriate formalizations of orientation state spaces, operators, and refusal/non-closure policies.


\nocite{*}
\bibliographystyle{unsrt}  
\bibliography{references}
\end{document}